\indent ಹೀಗೆ, ಕ್ರಿಯೆ ಸಂಭವಿಸುವಲ್ಲಿ ಸಂಶಯ ಅಥವಾ ಸಂದೇಹ ಪ್ರಕಟಿಸುವ ಕ್ರಿಯಾಪದಗಳಿಗೆ ಸಂಭಾವನಾರ್ಥಕ ಕ್ರಿಯಾಪದಗಳೆಂದು ಹೆಸರು. ಧಾತುಗಳ ಮೇಲೆ ಏಕವಚನದಲ್ಲಿ ಆನು, ಆಳು, ಈತು, ಈಯೆ, ಏನು ಎಂಬ ಆಖ್ಯಾತ ಪ್ರತ್ಯಯಗಳೂ ಬಹುವಚನದಲ್ಲಿ ಆರು, ಆವು, ಈರಿ, ಏವು—ಎಂಬ ಆಖ್ಯಾತ ಪ್ರತ್ಯಯಗಳೂ ಸೇರಿ ಸಂಭಾವನಾರ್ಥಕ ಕ್ರಿಯಾಪದ ಗಳಾಗುವುವು.\par

\begin{kbvrule}
\noindent\kbvbullet\hspace{0.5em}ಸಾಕು, ಬೇಕು, ಇಲ್ಲ, ಹೌದು, ಬೇಡ, ಅಲ್ಲ, ಉಂಟು ಇತ್ಯಾದಿ ಅವ್ಯಯಗಳನ್ನು ಕ್ರಿಯಾರ್ಥಕ ಪದಗಳೆನ್ನಬಹುದು. ಕ್ರಿಯಾಪದಗಳ ಅರ್ಥ ದಲ್ಲಿ ಬಳಸುವ ಈ ಅವ್ಯಯಗಳನ್ನು ಕ್ರಿಯಾರ್ಥಕಾವ್ಯಯಗಳೆನ್ನುವರು.
\end{kbvrule}

\begin{kbvderiv}
\vspace{4pt plus 3pt minus 1pt}
\noindent{\renewcommand{\arraystretch}{1.24}\setlength{\tabcolsep}{5pt}
\begin{tabularx}{\textwidth}{X X}
ಉದಾ : {[}೧{]} ಮಾತು ಸಾಕು, ಕೆಲಸ ಬೇಕು & {[}೨{]} ನನ್ನ ಜೇಬಿ \\
ನಲ್ಲಿ ಹಣ ಇಲ್ಲ & {[}೩{]} ನಿನ್ನ ಸಹವಾಸ ಬೇಡ. {[}೪{]} ಇದು ರೇಷ್ಮೆ ಸೀರೆ \\
\end{tabularx}}
\vspace{4pt plus 3pt minus 1pt}
\end{kbvderiv}

\indent ಅಲ್ಲ. (೫) ಅದು ಧ್ರುವ ನಕ್ಷತ್ರವೇ ಹೌದು. (೬) ಆಕೆಯಲ್ಲಿ ನನಗೆ ಪ್ರೀತಿ ಉಂಟು.\par

\kbvsubtitle{ಕರ್ತರಿ ಮತ್ತು ಕರ್ಮಣಿ ಪ್ರಯೋಗಗಳು}

\indent ವಾಕ್ಯಗಳ ಎರಡು ಶ್ರೇಣಿಗಳನ್ನು ಇಲ್ಲಿ ಕೊಟ್ಟಿದೆ. ಪ್ರತಿಯೊಂದು ಶ್ರೇಣಿಯಲ್ಲೂ ಮೂರು ವಾಕ್ಯಗಳಿವೆ, ಆ ವಾಕ್ಯಗಳನ್ನೂ ಅವುಗಳಲ್ಲಿರುವ ಕ್ರಿಯಾಪದಗಳನ್ನು ಗಮನಿಸಿರಿ.\par

\kbvsechead{ಮೊದಲನೆಯ ಶ್ರೇಣಿ :}

\indent ರಾಮನು ಹೊಲವನ್ನು ಮಾರುತ್ತಾನೆ. ...(ಮಾರು + ಉತ್ತ + ಆನೆ)\par

\indent ಸೀತೆಯು ಹೊಲವನ್ನು ಮಾರಿದಳು ....(ಮಾರು + ದ + ಅಳು)\par

\indent ಅವರು ಹೊಲವನ್ನು ಮಾರುವರು ....(ಮಾರು + ವ + ಅರು)\par

\kbvsechead{ಎರಡನೆಯ ಶ್ರೇಣಿ :}

\noindent\hspace*{2.2em}ರಾಮನಿಂದ ಹೊಲವು ಮಾರಲ್ಪಡುತ್ತದೆ.\par

\indent ...(ಮಾರು + ಅಲ್ಪಡು + ಉತ್ತ + ಅದೆ) ಸೀತೆಯಿಂದ ಹೊಲವು ಮಾರಲ್ಪಟ್ಟಿತು.\par

\indent ...(ಮಾರು + ಅಲ್ಪಡು + ದ + ಇತು)\par
